\documentclass{amsart}
\usepackage{geometry,amssymb,color,hyperref}
\usepackage{graphicx} % Required for inserting images
\usepackage{mathtools}

\theoremstyle{plain}
\newtheorem{theorem}{Theorem}

\theoremstyle{definition}
\newtheorem{definition}{Definition}

\theoremstyle{remark}
\newtheorem{remark}{Remark}
\newtheorem{example}{Example}

\newcommand{\IZ}{\mathbb Z}

\title{There exist Steiner systems $S(2,7,505)$}
\author{Ivan Hetman}
\subjclass{51E05, 51E10}
\date{April 2026}

\begin{document}
    \begin{abstract} In this note two Steiner systems $S(2,7,505)$ are presented. This resolves one of $21$ undecided cases for block designs with block length $7$, mentioned in Handbook of Combinatorial Designs.
    \end{abstract}

    \maketitle

    Given three integers $t$, $k$, $v$ such that $2 \le t < k < v$, a Steiner system $S(t, k, v)$ is a set $V$ of cardinality $v$ together with a family $\mathcal B$ of $k$-subsets of $V$ (called \textit{blocks}) such that every $t$-element subset of $V$ is contained in exactly one block.
    Handbook of Combinatorial Designs~\cite[II.3.5-6]{HoCD} lists $21$ numbers $v$ for which the existence of Steiner systems $S(2,7,v)$ is unknown.
    In this note we present two Steiner systems $S(2,7,505)$ thus resolving one of the above $21$ undecided cases.

    Difference families are convenient and compact way to present Steiner systems.
    Given a group $G$ of cardinality $v$, we look for a family $\mathcal B$ of $k$-element subsets (called \emph{base blocks}) of $G$ such that the pair $(G, \{Bg:B\in\mathcal B,g\in G\})$ is a Steiner system $S(2,k,v)$.
    For a base block $B$ denote by $\Delta B=\{ac^{-1}:a,c\in B, a\not= c\}$ its set of differences.
    It is well-known that when the group $G$ is cyclic $\IZ_v=\{0,\dots,v-1\}$ endowed with addition modulo $v$, and $v,k$ are coprime, then cardinality of all base blocks set of differences should be equal to $k(k-1)$ and all sets of differences should be disjoint.
    Using that knowledge, it is easy to verify that two families presented below indeed are difference families and they produce $S(2,7,505)$.
    Checking for isomorphism and automorphism group is more complicated and requires specialized software like \texttt{nauty}~\cite{nauty} or calculating fingerprint introduced in~\cite{Het}.
    Difference families will be preceded with design automorphism group order and fingerprint to ensure that corresponding designs are not isomorphic.

    \begin{enumerate}
    \item $2525\;\{1=15150, 2=424200, 3=14710650, 4=142803900, 5=475800900\} \newline \{\{0, 1, 3, 7, 47, 133, 284\}, \{0, 5, 70, 100, 173, 185, 476\}, \{0, 8, 82, 199, 248, 391, 474\}, \newline \{0, 9, 262, 298, 370, 386, 439\}, \{0, 10, 137, 156, 213, 233, 246\}, \{0, 11, 182, 250, 274, 414, 462\}, \newline \{0, 14, 78, 172, 366, 421, 477\}, \{0, 15, 136, 157, 174, 322, 344\}, \{0, 18, 181, 304, 330, 371, 442\}, \newline \{0, 23, 58, 120, 227, 287, 374\}, \{0, 25, 105, 150, 209, 260, 310\}, \{0, 27, 145, 206, 238, 416, 453\}\}$
    \item $2525\;\{2=444400, 3=13483500, 4=139198200, 5=480628700\} \newline \{\{0, 1, 3, 7, 119, 242, 341\}, \{0, 5, 51, 63, 95, 254, 287\}, \{0, 8, 261, 297, 369, 388, 417\}, \newline \{0, 9, 26, 85, 357, 428, 490\}, \{0, 10, 170, 310, 391, 455, 492\}, \{0, 11, 93, 113, 279, 384, 422\}, \newline \{0, 14, 161, 290, 368, 421, 477\}, \{0, 16, 89, 158, 273, 327, 453\}, \{0, 18, 45, 79, 217, 304, 371\}, \newline \{0, 21, 43, 149, 206, 317, 456\}, \{0, 25, 135, 175, 328, 375, 450\}, \{0, 30, 227, 258, 324, 431, 470\}\}$
    \end{enumerate}

    It is possible to obtain more designs using mirroring technique.

    \begin{definition}
        A \emph{mirror} of a family $\mathcal B=\{B_1,\dots,B_n\}$ is a family obtained by replacing some sets with their inverses.
        More formally, for any of the $2^n$ functions $\mu:\{1,\dots,n\}\to \{1, -1\}$, a family $\mathcal B^\mu = \{B_1^{\mu(1)},\dots,B_n^{\mu(n)}\}$ is a mirror of $\mathcal B$.
    \end{definition}

    Computer calculations show that from two presented difference families it is possible to obtain $832$ non-isomorphic designs, $16$ with automorphism group of order $2525$, and $816$ with automorphism group of order $505$.

    \begin{thebibliography}{}

        \bibitem{HoCD} {\em Handbook of Combinatorial Designs}. Edited by Charles J. Colbourn and Jeffrey H. Dinitz. Second edition. Discrete Mathematics and its Applications (Boca Raton). Chapman \& Hall/CRC, Boca Raton, FL, 2007.

        \bibitem{nauty} B.~D.~McKay, A.~Piperino, {\em Practical Graph Isomorphism, II}, Journal of Symbolic Computation, \textbf{60}, 94--112 (2014).

        \bibitem{Het} I.~Hetman, {\em There Exist Steiner Systems $S(2,8,225)$ and $S(2,9,289)$}, Journal of Combinatorial Designs (2026).
    \end{thebibliography}
\end{document}